% Options for packages loaded elsewhere
\PassOptionsToPackage{unicode}{hyperref}
\PassOptionsToPackage{hyphens}{url}
\documentclass[
  11pt,
]{article}
\usepackage{xcolor}
\usepackage[margin=1in]{geometry}
\usepackage{amsmath,amssymb}
\setcounter{secnumdepth}{-\maxdimen} % remove section numbering
\usepackage{iftex}
\ifPDFTeX
  \usepackage[T1]{fontenc}
  \usepackage[utf8]{inputenc}
  \usepackage{textcomp} % provide euro and other symbols
\else % if luatex or xetex
  \usepackage{unicode-math} % this also loads fontspec
  \defaultfontfeatures{Scale=MatchLowercase}
  \defaultfontfeatures[\rmfamily]{Ligatures=TeX,Scale=1}
\fi
\usepackage{lmodern}
\ifPDFTeX\else
  % xetex/luatex font selection
\fi
% Use upquote if available, for straight quotes in verbatim environments
\IfFileExists{upquote.sty}{\usepackage{upquote}}{}
\IfFileExists{microtype.sty}{% use microtype if available
  \usepackage[]{microtype}
  \UseMicrotypeSet[protrusion]{basicmath} % disable protrusion for tt fonts
}{}
\makeatletter
\@ifundefined{KOMAClassName}{% if non-KOMA class
  \IfFileExists{parskip.sty}{%
    \usepackage{parskip}
  }{% else
    \setlength{\parindent}{0pt}
    \setlength{\parskip}{6pt plus 2pt minus 1pt}}
}{% if KOMA class
  \KOMAoptions{parskip=half}}
\makeatother
\usepackage{longtable,booktabs,array}
\usepackage{caption}
\captionsetup[table]{skip=6pt}
\newcounter{none} % for unnumbered tables
\usepackage{calc} % for calculating minipage widths
% Correct order of tables after \paragraph or \subparagraph
\usepackage{etoolbox}
\makeatletter
\patchcmd\longtable{\par}{\if@noskipsec\mbox{}\fi\par}{}{}
\makeatother
% Allow footnotes in longtable head/foot
\IfFileExists{footnotehyper.sty}{\usepackage{footnotehyper}}{\usepackage{footnote}}
\makesavenoteenv{longtable}
\setlength{\emergencystretch}{3em} % prevent overfull lines
\providecommand{\tightlist}{%
  \setlength{\itemsep}{0pt}\setlength{\parskip}{0pt}}
\usepackage{bookmark}
\IfFileExists{xurl.sty}{\usepackage{xurl}}{} % add URL line breaks if available
\urlstyle{same}
% fallback for those not using the hyperref driver hyperxmp:
\makeatletter
\@ifundefined{xmpquote}{\newcommand{\xmpquote}[1]{#1}}{}
\makeatother
\hypersetup{
  hidelinks,
  pdfcreator={LaTeX via pandoc}}

\author{}
\date{}

\begin{document}

\section{RECYM: An Automated Multi-Feeder Workflow for Demand Allocation
and Load-Flow Studies in CYMDIST Distribution
Networks}\label{recym-an-automated-multi-feeder-workflow-for-demand-allocation-and-load-flow-studies-in-cymdist-distribution-networks}

\textbf{Draft manuscript — IEEE PES Transactions structure (full RECYM
scope)}\\
\textbf{Target venue:} \emph{IEEE Transactions on Power Delivery} (alt.:
\emph{IEEE Trans. Smart Grid} / \emph{IEEE Access})\\
\textbf{Citation style:} IEEE numbered \texttt{{[}n{]}} in order of
appearance\\
\textbf{Status:} Draft v0.2 — whole project (SPA §§1–7, API/jobs,
pipeline, CymPy/COM, multi-feeder, PA217 evidence)\\
\textbf{Authors / affiliation:} \emph{{[}To be completed{]}}\\
\textbf{Corresponding author:} \emph{{[}email{]}}

\subsubsection{IEEE / PES numerals enforced in this
draft}\label{ieee-pes-numerals-enforced-in-this-draft}

{\def\LTcaptype{none} % do not increment counter
\begin{longtable}[]{@{}
  >{\raggedright\arraybackslash}p{(\linewidth - 4\tabcolsep) * \real{0.4062}}
  >{\raggedright\arraybackslash}p{(\linewidth - 4\tabcolsep) * \real{0.2188}}
  >{\raggedright\arraybackslash}p{(\linewidth - 4\tabcolsep) * \real{0.3750}}@{}}
\toprule\noalign{}
\begin{minipage}[b]{\linewidth}\raggedright
Requirement
\end{minipage} & \begin{minipage}[b]{\linewidth}\raggedright
Limit
\end{minipage} & \begin{minipage}[b]{\linewidth}\raggedright
This draft
\end{minipage} \\
\midrule\noalign{}
\endhead
\bottomrule\noalign{}
\endlastfoot
Abstract & \textbf{150–250 words}, one paragraph; no citations/equations
& \textasciitilde159 words \\
Index Terms & Alphabetical keywords & Yes (separate block) \\
Body sections & Roman \textbf{I–VII} (+ A/B subsections) & Yes \\
First PES submission & \textbf{≤ 10 pages} (2-column, figs/refs/bios
included) & Keep when compiling IEEEtran \\
Citations & \texttt{{[}1{]}}, \texttt{{[}2{]}–{[}4{]}} & Yes \\
Units & SI (kW, kvar, kWh, kV, p.u.) & Yes \\
\end{longtable}
}

\begin{quote}
Source artifacts: \texttt{docs/ARQUITECTURA.md},
\texttt{docs/FLUJO\_TRABAJO.md},\\
\texttt{data/output/feeders/PA217/} (session, allocation, cierre\_1\_7,
inventory, images).\\
Copy into official IEEEtran / IEEE Word template; trim Appendix if PDF
\textgreater{} 10 pages.
\end{quote}

\begin{center}\rule{0.5\linewidth}{0.5pt}\end{center}

\subsection{Abstract}\label{abstract}

Distribution utilities need repeatable procedures to clean network
models, allocate measured substation demand, inject important-customer
and prospective spot loads, and compare situational versus projected
operating cases. This paper presents RECYM, a multi-feeder software
suite that couples CYMDIST/CymPy with a seven-step single-page
application and a batch pipeline. The workflow integrates network
diagnostics, large-customer energy/power mapping to secondary
distribution substations, kilowatt-hour-based load allocation via the
CYMDIST COM API, locked three-phase spot loads, dual load-flow
scenarios, and automated delivery reports. On feeder PA217 (22.9 kV,
Electro Dunas), the campaign closed with zero diagnostic problems,
important customers verified within configured energy/power tolerances,
headroom-balanced allocation (9537.9 kW versus 9537.88 kW), a 1420 kW
prospective spot load, and successful situational and projected load
flows. The same architecture scales to a catalog of about 96 feeders
sharing one Access database and per-feeder study files.

\subsection{Index Terms}\label{index-terms}

CYMDIST, demand studies, distribution networks, load allocation, load
flow analysis, power system modeling, software tools.

\begin{center}\rule{0.5\linewidth}{0.5pt}\end{center}

\subsection{I. Introduction}\label{i.-introduction}

Modern distribution planning relies on detailed feeder models to assess
voltage profiles, thermal loading, and the impact of new customers
{[}1{]}, {[}2{]}. Commercial packages such as CYMDIST provide load
allocation and unbalanced load-flow engines used widely by utilities
{[}3{]}. In practice, however, campaigns still depend on manual GUI
steps: correcting DEFAULT equipment, transferring billing energy and
demand to SED nodes, locking loads, connecting prospective spot loads,
and regenerating reports for situational (existing) versus projected
(with new load) cases.

Incomplete or inconsistent models degrade allocation and load-flow
quality {[}1{]}, {[}4{]}. Open tooling such as OpenDSS {[}5{]} and
pandapower {[}6{]} supports research automation, yet many Latin American
utilities already standardize on CYMDIST study (\texttt{.zxst}) and
Access (\texttt{.mdb}) assets. Bridging enterprise models with scripted,
auditable workflows remains a gap {[}7{]}, {[}8{]}.

This paper introduces \textbf{RECYM} (\emph{Redes / Electro Dunas
CYMDIST Management}), a suite that:

\begin{enumerate}
\def\labelenumi{\arabic{enumi}.}
\tightlist
\item
  Provides a layered architecture (SPA UI, REST/job API, Python
  pipeline, CymPy/COM core) for multi-feeder operation;
\item
  Encodes a seven-step campaign from context and model-quality gate
  through dual load flows and delivery reports;
\item
  Demonstrates end-to-end closure on a real 22.9\,kV feeder (PA217) with
  quantitative precision and balance metrics.
\end{enumerate}

Section II reviews related work. Section III describes the architecture.
Section IV details the proposed workflow. Section V presents the case
study. Section VI discusses results. Section VII concludes.

\begin{center}\rule{0.5\linewidth}{0.5pt}\end{center}

\subsection{II. Related Work}\label{ii.-related-work}

\subsubsection{A. Distribution Modeling and Load
Allocation}\label{a.-distribution-modeling-and-load-allocation}

Kersting {[}1{]} formalizes three-phase distribution modeling and load
representation. Load allocation methods distribute measured headroom
demand across customer meters using energy, demand, or transformer kVA
as weighting factors {[}3{]}, {[}9{]}. CYMDIST implements
Connected+Total allocation templates with consumption (kWh), diversity,
and loss factors {[}3{]}. Accuracy depends on locked large customers and
clean connectivity {[}9{]}, {[}10{]}.

\subsubsection{B. Automation of Utility
Studies}\label{b.-automation-of-utility-studies}

Dugan and McDermott {[}5{]} popularized scriptable distribution
analysis. Python ecosystems (pandapower, PyPSA) emphasize open data
models {[}6{]}, {[}11{]}. COM/API automation of commercial tools is less
documented in the literature but is operationally critical when the
system of record is proprietary {[}7{]}, {[}8{]}. RECYM targets that
niche: reproducible CYMDIST campaigns without abandoning the utility’s
native database.

\subsubsection{C. Capacity and New-Load
Studies}\label{c.-capacity-and-new-load-studies}

Hosting-capacity and interconnection studies compare base and
incremental cases under voltage and thermal limits {[}12{]}, {[}13{]}.
Regulatory interconnection frameworks (e.g., IEEE Std 1547 {[}14{]})
motivate clear situational vs.~projected narratives. RECYM
operationalizes that comparison at feeder scale for a new medium-voltage
customer request.

\begin{center}\rule{0.5\linewidth}{0.5pt}\end{center}

\subsection{III. System Architecture}\label{iii.-system-architecture}

\subsubsection{A. Layered Design}\label{a.-layered-design}

Fig. 1 (see \texttt{docs/ARQUITECTURA.md}) summarizes five layers:

{\def\LTcaptype{none} % do not increment counter
\begin{longtable}[]{@{}
  >{\raggedright\arraybackslash}p{(\linewidth - 2\tabcolsep) * \real{0.5385}}
  >{\raggedright\arraybackslash}p{(\linewidth - 2\tabcolsep) * \real{0.4615}}@{}}
\toprule\noalign{}
\begin{minipage}[b]{\linewidth}\raggedright
Layer
\end{minipage} & \begin{minipage}[b]{\linewidth}\raggedright
Role
\end{minipage} \\
\midrule\noalign{}
\endhead
\bottomrule\noalign{}
\endlastfoot
Presentation & React SPA (§§1–7), Vite build served at port 5055 \\
API & FastAPI jobs/SSE + Flask handler bridge for \texttt{/api/*} \\
Domain & \texttt{pipeline/}, \texttt{analysis/},
\texttt{optimization/} \\
Core & \texttt{cympy\_adapter}, \texttt{cymdist\_com},
\texttt{feeder\_context}, Excel I/O \\
External data & CYMDIST \texttt{.zxst} studies, shared \texttt{.mdb},
Excel workbooks \\
\end{longtable}
}

Long-running actions (diagnostics, allocation, load flow) run as
asynchronous jobs with Server-Sent Events (SSE) progress, avoiding UI
timeouts during COM calls.

\subsubsection{B. Multi-Feeder
Configuration}\label{b.-multi-feeder-configuration}

Global settings (\texttt{config/settings.json}) define CYME paths, the
shared Access database, default equipment IDs, precision tolerances, and
the batch \texttt{run\_sequence}. Each feeder file
(\texttt{config/feeders/\textless{}ID\textgreater{}.json}) binds
\texttt{network\_id}, study path, control/catalog workbooks, and output
directory. Selection follows CLI (\texttt{-\/-feeder}), environment
(\texttt{RECYM\_FEEDER}), UI header \texttt{X-Feeder}, or
\texttt{active\_feeder}. The catalog used in this work contains on the
order of \textbf{96} Electro Dunas feeders plus study ELD for
system-wide diagnostics.

\subsubsection{C. CYMDIST Integration}\label{c.-cymdist-integration}

Write operations require a usable study file and
\texttt{dry\_run=false}. Load allocation and load flow preferentially
use the \textbf{COM} engine (\texttt{loadflow\_engine:\ COM}) for
reliability when native CymPy allocation is unavailable. Customer fields
mapped by RECYM include KWH consumption, kW/kvar,
\texttt{ConnectionStatus} (Connected/Disconnected), and
\texttt{LockDuringLoadAllocation} (Locked/Unlocked) {[}3{]}.

\begin{center}\rule{0.5\linewidth}{0.5pt}\end{center}

\subsection{IV. Proposed Workflow}\label{iv.-proposed-workflow}

The SPA enforces a numbered campaign (Fig. 2 /
\texttt{docs/FLUJO\_TRABAJO.md}).

\subsubsection{A. §1 Context and
Headroom}\label{a.-1-context-and-headroom}

The operator selects database and study, then loads feeder headroom
measurement (P, Q, optionally from metering Excel). CYMDIST network
allocation templates are set to Total ON, downstream consumption in kWh,
load factor (default 65\,\%), and loss factor (k=0.3).

\subsubsection{B. §2 Model-Quality Gate}\label{b.-2-model-quality-gate}

NetworkDiagnostic runs at feeder or system (96-network / ELD) scope.
Corrections for DEFAULT equipment and base voltages are proposed and
applied. The gate requires zero Error/Warning/Hint messages and
load-flow convergence before demand steps.

\subsubsection{C. §3 Important Customers and
Allocation}\label{c.-3-important-customers-and-allocation}

Billing supply and “important customers” tables are crossed by NIS. For
each included SED:

\begin{itemize}
\tightlist
\item
  Energy \textbf{EA} → CYMDIST \textbf{KWH} (consumption);
\item
  Demand \textbf{Pot} → \textbf{kW}, Locked during allocation;
\item
  Excluded rows → Disconnected + 0\,kW.
\end{itemize}

Load allocation then distributes residual headroom by kWh weights.
Precision checks enforce (\textbar{}\Delta\mathrm{kWh}\textbar{}\le 0.5)
and (\textbar{}\Delta\mathrm{kW}\textbar{}\le 0.01).

\subsubsection{D. §4 Prospective Spot
Load}\label{d.-4-prospective-spot-load}

A concentrated three-phase spot load is attached to an \emph{existing}
node (no new topology). Active and reactive power are split equally:

{[} P\_A=P\_B=P\_C=\frac{P_{3\phi}}{3},\qquad
Q\_A=Q\_B=Q\_C=\frac{Q_{3\phi}}{3}. {]}

The load is Locked. After §4, allocation (§3.3) is not re-run; only load
flow follows.

\subsubsection{E. §5 Dual Load-Flow
Scenarios}\label{e.-5-dual-load-flow-scenarios}

{\def\LTcaptype{none} % do not increment counter
\begin{longtable}[]{@{}lll@{}}
\toprule\noalign{}
Scenario & Spot-load §4 & Purpose \\
\midrule\noalign{}
\endhead
\bottomrule\noalign{}
\endlastfoot
Situational & Disconnected & Existing network condition \\
Projected & Connected & With prospective customer \\
General & Unchanged & Ad-hoc LF \\
\end{longtable}
}

\subsubsection{F. §§6–7 Reports and
Suite}\label{f.-67-reports-and-suite}

§6 fills delivery Word/PDF reports from OCR meta (customer, kW request)
and scenario images. §7 exposes optional device optimization (reclosers,
regulators, capacitors) and batch suite tools.

\textbf{Algorithm 1} (scenario load flow):

\begin{verbatim}
Input: study S, spot-load set L, scenario ∈ {situational, projected}
if scenario = situational: disconnect each ℓ ∈ L
if scenario = projected:   connect each ℓ ∈ L
Run COM LoadFlow on S
Export voltage/loading summaries and map images
\end{verbatim}

\begin{center}\rule{0.5\linewidth}{0.5pt}\end{center}

\subsection{V. Case Study}\label{v.-case-study}

\subsubsection{A. Network and Software}\label{a.-network-and-software}

\begin{itemize}
\tightlist
\item
  \textbf{Utility:} Electro Dunas (Peru).
\item
  \textbf{Feeder:} PA217, network \texttt{NET\_2030\_179\_PA217},
  (V\_\{LL\}=22.9,\mathrm{kV}), SET Paracas.
\item
  \textbf{Study / DB:} \texttt{PA217.zxst}, Access
  \texttt{20260919.mdb}.
\item
  \textbf{Software:} CYMDIST 9.2, CymPy, Python 3.7 win32, RECYM SPA v6.
\item
  \textbf{Prospective customer:} San Fernando S.A., incubadora Sur
  project, request \textbf{1420\,kW} (feasibility at MV).
\end{itemize}

\subsubsection{B. Headroom Measurement}\label{b.-headroom-measurement}

From metering session (\texttt{L-PA235}, Excel \emph{SISTEMA Pisco.xls},
2026-03-03):

{\def\LTcaptype{none} % do not increment counter
\begin{longtable}[]{@{}lr@{}}
\toprule\noalign{}
Quantity & Value \\
\midrule\noalign{}
\endhead
\bottomrule\noalign{}
\endlastfoot
(P) & 9537.88\,kW \\
(Q) & 2587.65\,kvar \\
(S) & 9882.67\,kVA \\
Avg. (P) & 6204.65\,kW \\
Load factor & 65.05\,\% \\
\end{longtable}
}

\subsubsection{C. Metrics}\label{c.-metrics}

Diagnostic message counts; customer EA/Pot precision; allocation
headroom balance; spot-load phase split; situational/projected LF
success; end-to-end campaign closure
(\texttt{cierre\_1\_7\_report.json}).

\begin{center}\rule{0.5\linewidth}{0.5pt}\end{center}

\subsection{VI. Results and
Discussion}\label{vi.-results-and-discussion}

\subsubsection{A. Model Quality}\label{a.-model-quality}

Post-campaign diagnostics for PA217 reported \textbf{0} messages and
\textbf{0} problems (before/after summaries empty), satisfying the §2
gate for the validated run.

\subsubsection{B. Important-Customer
Precision}\label{b.-important-customer-precision}

Twelve candidate rows were built; \textbf{11} SED mappings passed
precision tests ((n\_\{\mathrm{fail}\}=0)); one supply without SED was
skipped. Representative verified customers include Caliza Cementos Inca
(3800.57\,kW / 1.75\,GWh), Minerales Paracas (1273.76\,kW), and Emsalsa
(458.00\,kW). Absolute errors were within configured tolerances (Table
I).

\textbf{TABLE I}\\
IMPORTANT-CUSTOMER PRECISION ON PA217 (SAMPLE)

{\def\LTcaptype{none} % do not increment counter
\begin{longtable}[]{@{}
  >{\raggedright\arraybackslash}p{(\linewidth - 12\tabcolsep) * \real{0.1389}}
  >{\raggedright\arraybackslash}p{(\linewidth - 12\tabcolsep) * \real{0.0694}}
  >{\raggedright\arraybackslash}p{(\linewidth - 12\tabcolsep) * \real{0.1389}}
  >{\raggedright\arraybackslash}p{(\linewidth - 12\tabcolsep) * \real{0.1389}}
  >{\raggedright\arraybackslash}p{(\linewidth - 12\tabcolsep) * \real{0.2083}}
  >{\raggedright\arraybackslash}p{(\linewidth - 12\tabcolsep) * \real{0.1944}}
  >{\raggedright\arraybackslash}p{(\linewidth - 12\tabcolsep) * \real{0.1111}}@{}}
\toprule\noalign{}
\begin{minipage}[b]{\linewidth}\raggedright
Customer
\end{minipage} & \begin{minipage}[b]{\linewidth}\raggedright
SED
\end{minipage} & \begin{minipage}[b]{\linewidth}\raggedright
EA (kWh)
\end{minipage} & \begin{minipage}[b]{\linewidth}\raggedright
Pot (kW)
\end{minipage} & \begin{minipage}[b]{\linewidth}\raggedright
(\Delta)kWh
\end{minipage} & \begin{minipage}[b]{\linewidth}\raggedright
(\Delta)kW
\end{minipage} & \begin{minipage}[b]{\linewidth}\raggedright
Status
\end{minipage} \\
\midrule\noalign{}
\endhead
\bottomrule\noalign{}
\endlastfoot
Caliza cementos inca & SE30940 & 1\,747\,551.27 & 3800.57 & 0.0 & 0.0 &
OK \\
Minerales Paracas & SE30480 & 678\,852.12 & 1273.76 & 0.0 & 0.0 & OK \\
LJMetales & SE30642 & 527\,079.81 & 956.49 & 0.0 & 0.0 & OK \\
Emsalsa & SE30353 & 157\,621.44 & 458.00 & 0.0 & 0.0 & OK \\
Uvica & SE30424 & 167\,101.44 & 366.30 & 0.0 & 0.0 & OK \\
\end{longtable}
}

\subsubsection{C. Load Allocation
Balance}\label{c.-load-allocation-balance}

COM kWh-based LoadAllocation yielded (\sum P = 9537.9,\mathrm{kW})
against headroom (P=9537.88,\mathrm{kW}) ((\Delta\approx 0\%)), with
validation mode reporting balance OK over 257 modeled loads.

\subsubsection{D. Prospective Spot Load}\label{d.-prospective-spot-load}

A locked spot load of \textbf{1420\,kW} was placed at node
\texttt{NODE\_1080\_320357\_MI-17145\_2} (UTM zone 18S; approx.
13.862°\,S, 76.132°\,W). An independent integral check on device
\texttt{SAN\_FERNANDO1} confirmed phase powers
(466.67/466.67/466.67,\mathrm{kW}) for a 1400\,kW three-phase injection
(equal split), and the orphan zero-kW device was disconnected.

\subsubsection{E. Dual Load Flows and Campaign
Closure}\label{e.-dual-load-flows-and-campaign-closure}

Situational and projected load flows both completed successfully. The
automated closure report recorded \textbf{21/21} steps OK
(\texttt{closed:\ true}), including context, diagnostics, EA/Pot apply,
allocation, both LF scenarios, informe assembly, and suite checks.

\subsubsection{F. Discussion and
Limitations}\label{f.-discussion-and-limitations}

RECYM reduces GUI chatter and produces auditable JSON/CSV artifacts per
feeder. Limitations include: (i) Windows/COM session coupling to
CYMDIST; (ii) dependence on study-file integrity; (iii) native CymPy
allocation may be unavailable, requiring COM fallback; (iv) optimization
(§7) is exposed but not the focus of the PA217 closure metrics. Results
are for one feeder; multi-feeder statistical aggregation is future work.

\begin{center}\rule{0.5\linewidth}{0.5pt}\end{center}

\subsection{VII. Conclusion}\label{vii.-conclusion}

This paper presented RECYM, an automated multi-feeder workflow for
CYMDIST demand and load-flow studies. The seven-step SPA and batch
pipeline combine model-quality gating, important-customer injection, kWh
load allocation, locked spot loads, and situational/projected load
flows—covering the full product stack (UI, API/jobs, domain pipelines,
CymPy/COM, multi-feeder config, and report generation). On Electro Dunas
feeder PA217, the method achieved zero diagnostic issues, verified
customer precision, headroom-balanced allocation, dual load-flow closure
for a 1420 kW prospective MV customer, and a 21/21 automated campaign
closure. Future work includes fleet-wide KPI dashboards across
\textasciitilde96 feeders, tighter regulatory voltage/thermal scoring,
and production hardening of §7 device optimization.

\begin{center}\rule{0.5\linewidth}{0.5pt}\end{center}

\subsection{\texorpdfstring{Appendix A — Batch \texttt{run\_sequence}
(excerpt)}{Appendix A — Batch run\_sequence (excerpt)}}\label{appendix-a-batch-run_sequence-excerpt}

\texttt{validate\_inputs} → \texttt{network\_diagnostic} →
\texttt{sync\_equipment} → \texttt{build\_corrections} →
\texttt{bulk\_fix} → \texttt{fix\_base\_voltages} →
\texttt{network\_diagnostic\_after} → \texttt{inventory\_loads} →
\texttt{build\_clientes} → \texttt{apply\_clientes} →
\texttt{verify\_precision} → \texttt{demand\_allocation} →
\texttt{report}.

\subsection{Appendix B — Figure
Sources}\label{appendix-b-figure-sources}

{\def\LTcaptype{none} % do not increment counter
\begin{longtable}[]{@{}
  >{\raggedright\arraybackslash}p{(\linewidth - 2\tabcolsep) * \real{0.2727}}
  >{\raggedright\arraybackslash}p{(\linewidth - 2\tabcolsep) * \real{0.7273}}@{}}
\toprule\noalign{}
\begin{minipage}[b]{\linewidth}\raggedright
Fig.
\end{minipage} & \begin{minipage}[b]{\linewidth}\raggedright
Suggested file
\end{minipage} \\
\midrule\noalign{}
\endhead
\bottomrule\noalign{}
\endlastfoot
1 Architecture & Diagram in \texttt{docs/ARQUITECTURA.md} \\
2 Workflow & Mermaid in \texttt{docs/FLUJO\_TRABAJO.md} \\
3 Location map &
\texttt{data/output/feeders/PA217/informe\_images/topologia.png} \\
4–5 Voltage/loading & \texttt{situacional\_*.png},
\texttt{proyectado\_*.png} \\
\end{longtable}
}

\begin{center}\rule{0.5\linewidth}{0.5pt}\end{center}

\subsection{References}\label{references}

{[}1{]} W. H. Kersting, \emph{Distribution System Modeling and
Analysis}, 4th ed.~Boca Raton, FL, USA: CRC Press, 2017.

{[}2{]} T. Gönen, \emph{Electric Power Distribution Engineering}, 3rd
ed.~Boca Raton, FL, USA: CRC Press, 2014.

{[}3{]} Eaton CYME, \emph{CYMDIST Reference Manual / How-to Tutorials}
(Load Allocation, Load Flow), version 9.2. Saint-Bruno, QC, Canada:
Eaton, 2022.

{[}4{]} R. C. Dugan, M. F. McGranaghan, S. Santoso, and H. W. Beaty,
\emph{Electrical Power Systems Quality}, 3rd ed.~New York, NY, USA:
McGraw-Hill, 2012.

{[}5{]} R. C. Dugan and T. E. McDermott, “An open source platform for
collaborating on smart grid research,” in \emph{Proc. IEEE Power Energy
Soc. Gen.~Meeting}, Detroit, MI, USA, 2011, pp.~1–7.

{[}6{]} L. Thurner \emph{et al.}, “pandapower—An open-source Python tool
for convenient modeling, analysis, and optimization of electric power
systems,” \emph{IEEE Trans. Power Syst.}, vol.~33, no. 6, pp.~6510–6521,
Nov.~2018.

{[}7{]} J. Peppanen, M. J. Reno, M. Thakkar, S. Grijalva, and R. G.
Harley, “Leveraging AMI data for distribution system model calibration
and situational awareness,” \emph{IEEE Trans. Smart Grid}, vol.~6, no.
4, pp.~2000–2009, Jul.~2015.

{[}8{]} Y. Wang, Q. Chen, T. Hong, and C. Kang, “Review of smart meter
data analytics: Applications, methodologies, and challenges,” \emph{IEEE
Trans. Smart Grid}, vol.~10, no. 3, pp.~3125–3148, May 2019.

{[}9{]} A. K. Ghosh, D. L. Lubkeman, and R. H. Jones, “Load modeling for
distribution circuit state estimation,” \emph{IEEE Trans. Power Del.},
vol.~12, no. 2, pp.~999–1005, Apr.~1997.

{[}10{]} M. E. Baran and A. W. Kelley, “State estimation for real-time
monitoring of distribution systems,” \emph{IEEE Trans. Power Syst.},
vol.~9, no. 3, pp.~1601–1609, Aug.~1994.

{[}11{]} T. Brown, J. Hörsch, and D. Schlachtberger, “PyPSA: Python for
power system analysis,” \emph{J. Open Res. Softw.}, vol.~6, no. 1, Art.
no. 4, 2018.

{[}12{]} M. Rylander, J. Smith, and W. Sunderman, “Streamlined method
for determining distribution system hosting capacity,” in \emph{Proc.
IEEE Rural Electric Power Conf.}, Asheville, NC, USA, 2015, pp.~3–9.

{[}13{]} F. Ding and B. Mather, “On distributed PV hosting capacity
estimation, sensitivity study, and improvement,” \emph{IEEE Trans.
Sustain. Energy}, vol.~8, no. 3, pp.~1010–1020, Jul.~2017.

{[}14{]} \emph{IEEE Standard for Interconnection and Interoperability of
Distributed Energy Resources with Associated Electric Power Systems
Interfaces}, IEEE Std 1547-2018, Apr.~2018.

{[}15{]} H. L. Willis, \emph{Power Distribution Planning Reference
Book}, 2nd ed.~New York, NY, USA: Marcel Dekker, 2004.

{[}16{]} J. A. Momoh, \emph{Electric Power Distribution, Automation,
Protection, and Control}. Boca Raton, FL, USA: CRC Press, 2008.

{[}17{]} A. Abur and A. G. Expósito, \emph{Power System State
Estimation: Theory and Implementation}. New York, NY, USA: Marcel
Dekker, 2004.

{[}18{]} OSINERGMIN, \emph{Normativa técnica de calidad de servicio
eléctrico / distribución} (Perú). Lima, Peru: Organismo Supervisor de la
Inversión en Energía y Minería. {[}Online{]}. Available:
https://www.osinergmin.gob.pe

{[}19{]} IEEE PES, “IEEE guide for electric power distribution
reliability indices,” IEEE Std 1366-2022, 2022.

{[}20{]} S. Conti, S. Raiti, G. Tina, and U. Vagliasindi, “Study of the
impact of PV generation on voltage profile in LV distribution networks,”
in \emph{Proc. IEEE Power Tech}, St.~Petersburg, Russia, 2005, pp.~1–6.

\end{document}
